
\section{Limitations}
\paragraph{Robot and compute resources are underutilized}

Coding agents do not fully utilize robot resource when they are reading logs, writing code, debugging, or waiting for the language-model backbone. Moreover, as we scale the number of robots, MRU decreases, while GPU active utilization increases (\Cref{fig:utilization}a). Compared to a single-robot setup, the composition of agent activity causes agent teams to spend more of their time summarizing peer branches and less time actually operating the robot. Coding agents may also fail to launch parallel training sessions to exhaust GPU resources.

\paragraph{Token cost grows super-linearly with fleet size.}
As the fleet size increases, token usage grows faster than the ideal linear trend. MTU remains close to the linear projection up to four agents, but rises sharply at eight agents (\Cref{fig:utilization}b). The total token budget required to obtain a successful policy shows the same trend, increasing much more rapidly than the corresponding reduction in wall-clock time (\Cref{fig:utilization}c). Thus, increasing the fleet size trades token efficiency for faster policy improvement: larger fleets reach success sooner, but require a disproportionately higher token budget.

% Finally, the hypothesis search uses a basic tree-expansion strategy with heuristic node selection; the persistent memory retrieves prior artifacts by filename and recency rather than by semantic similarity; offline component verification relies on cached sensor logs that the agent must curate manually. More sophisticated search (bandit-based selection or learned value functions over hypotheses), retrieval (embedding- or graph-structured memory of prior experiments), and verification (agent-authored unit tests over the toolkit) are all likely to substantially improve data efficiency.

% Second, the manipulation tasks we study are long-horizon as a whole, but on-robot RL is invoked only on the precise, contact-rich segment where classical primitives fail. The agent-authored task-focusing reset absorbs the rest: for pin insertion, the reset grasps and transports the pin using grasp synthesis \citep{anygrasp} and motion planning \citep{curobo}, hovering it above the insertion site so that RL is responsible only for the final precision-insertion step. This decomposition is exactly what makes real-world RL tractable under our hardware budget, but it presupposes a clean separation between a plannable prefix and a learnable bottleneck, which does not hold for long-horizon multi-stage manipulation, deformable-object manipulation, or settings in which the bottleneck shifts across episodes. Extending \method to agent-guided end-to-end RL on such tasks---where the agent must additionally reason about exploration, credit assignment, and curriculum over the full horizon---is an important direction we leave to future work. 

% \paragraph{Where the gap lies.}
% Across every frontier coding agent we evaluate, robots stand idle for the larger fraction of any research session, and the inefficiency compounds as the fleet grows. This is not a limit of the underlying language models, which continue to scale on the digital benchmarks they were trained for; it is a limit of the \emph{harness} that surrounds them. Current coding-agent loops are serial---read, think, then act---and lack the scheduling primitives needed to interleave language-model inference with physical experimentation. Closing this gap through asynchronous job queues, learned scheduling that overlaps thinking and acting, or summarized peer state that compresses per-iteration context is, in our view, the most immediate direction for follow-up work. The MRU and MCU metrics we propose are intended precisely to make this gap measurable and to allow direct comparison across systems as the harness improves.


